\section{The sharp no-waste rank bound}\label{sec:nowaste-sharp}

This section answers the final problem in the source manuscript.
The constant three can be replaced by one for simple graphs.
The proof uses the manuscript's elementary Erd\H{o}s--P\'osa lemma.
It does not use the asymptotic forest estimate.

\begin{theorem}\label{thm:nowaste-one}
Let $G$ be a connected simple graph, and let $W\subseteq G$ contain a cycle.
Suppose every outside part of an even subgraph is the outside part of a cycle.
Then
\[
 \beta(G)-\beta(W)\le |V(W)|+1.
\]
The additive constant one is sharp.
\end{theorem}

We first give the small classification needed in the proof.
A cyclic block means a maximal two-connected subgraph that contains a cycle.
Tree attachments do not change the cycle space.

\begin{lemma}\label{lem:nowaste-small-cores}
Suppose every nonempty even subgraph of a simple graph $F$ is a cycle.
Then $\beta(F)\le4$, and at most one component contains a cycle.
If $\beta(F)=2$, its unique cyclic block consists of three internally disjoint paths between two branch vertices.
If $\beta(F)=3$, that block is a subdivision of $K_4$.
If $\beta(F)=4$, that block is a subdivision of $K_{3,3}$.
All other edges in the cyclic component belong to trees attached to that block.
\end{lemma}

\begin{proof}
Two edge-disjoint cycles would have a union that is not a cycle.
Thus there is at most one cyclic block in the entire graph.
Every other edge is a bridge.
Remove the tree attachments and call the remaining cyclic block $B$.
It is two-connected and has minimum degree at least two.

No vertex of $B$ can have degree at least four.
Suppose $v$ did, and choose four edges incident to $v$.
Since $B-v$ is connected, it has an edge set whose odd-degree vertices are the four chosen neighbors.
Such an edge set exists by taking the appropriate join in a spanning tree.
Its union with the four chosen edges is even, with degree four at $v$.
This contradicts the hypothesis.
Hence the maximum degree of $B$ is three.

If every degree is two, then $B$ is a cycle and has rank one.
Otherwise suppress all degree-two vertices to obtain a connected cubic multigraph $J$.
Suppression preserves rank and the property concerning even subgraphs.
The manuscript's Erd\H{o}s--P\'osa lemma gives
$|E(J)|\le|V(J)|+3$.
As $J$ is cubic, it follows that $|V(J)|\le6$ and $\beta(J)\le4$.

For rank two, the cubic graph has two vertices.
It consists of three parallel edges, since the cyclic block has no bridge.
Its subdivision is the stated theta graph.
For rank three, the cubic graph has four vertices.
For rank four, it has six vertices.
In either case it has no loop or parallel pair.
Deleting such a cycle would leave at least as many edges as vertices, and therefore another edge-disjoint cycle.
The only simple cubic graph on four vertices is $K_4$.

In the six-vertex case, a triangle would also leave six edges after deletion, which again forces another cycle.
Thus $J$ is triangle-free.
Fix a vertex and its three neighbors.
The neighbors are independent and must each join both remaining vertices to have degree three.
This is exactly $K_{3,3}$.
\end{proof}

\begin{lemma}\label{lem:nowaste-path-obstruction}
Let $F=G-E(W)$, on the full vertex set of $G$.
Under the hypotheses of Theorem~\ref{thm:nowaste-one}, the following statements hold.
Every nonempty even subgraph of $F$ is a cycle.
Moreover, if $C$ is a cycle of $F$, then each component of $F-E(C)$ contains at most one vertex from each component of $W$.
\end{lemma}

\begin{proof}
An even edge set of $F$ is an outside part.
If a cycle of $G$ has this nonempty outside part, it contains that even edge set.
A cycle cannot contain a proper nonempty even edge set.
This proves the first assertion.

For the second assertion, suppose distinct $u,v$ lie in the same component of $W$ and the same component of $F-E(C)$.
Choose a $u$--$v$ path $P$ in $F-E(C)$ and a $u$--$v$ path $Q$ in $W$.
Their edge sets are disjoint.
The set $P\cup Q\cup C$ is even, since $P,Q$ have the same odd-degree vertices and $C$ is even.
Its outside part is $P\cup C$.
This contains the entire cycle $C$ and at least one further edge.
No cycle can have such an outside part, a contradiction.
\end{proof}

\begin{proof}[Proof of Theorem~\ref{thm:nowaste-one}]
Put $s=\operatorname{comp}(F)$, $c=\operatorname{comp}(W)$, $m=|V(W)|$, and $k=\beta(F)$.
Counting edges and vertices gives
\begin{equation}\label{eq:nowaste-rank-identity}
 \beta(G)-\beta(W)=m+k-s-c+1.
\end{equation}
It suffices to show $k\le s+c$.
Here $s,c\ge1$, so the assertion is immediate for $k\le2$.
By Lemmas~\ref{lem:nowaste-small-cores} and~\ref{lem:nowaste-path-obstruction}, only $k=3,4$ remain.

Choose a component $W_0$ of $W$ that contains a cycle.
It has at least three vertices, since $G$ is simple.
Each tree component of $F$ contains at most one vertex of $W_0$.
Indeed, the path between two such vertices would avoid any cycle in the cyclic component of $F$.

Suppose first that $k=3$.
The cyclic component of $F$ has a core that is a subdivision of $K_4$.
For each of the four triangles of $K_4$, delete the edges of its subdivided cycle.
One component of the remaining graph contains all four branch vertices.
Its core edges form the complementary spanning star.
Call these four components $T_1,\ldots,T_4$, including their surviving tree attachments.
Each $T_i$ contains at most one vertex of $W_0$, by Lemma~\ref{lem:nowaste-path-obstruction}.

Every vertex of the cyclic component belongs to at least two of these four components.
A branch vertex belongs to all four.
An internal vertex of a subdivided edge belongs to the two complementary stars that contain that edge.
A vertex in an attached tree has the same memberships as its attachment vertex.
Therefore the cyclic component contains at most two vertices of $W_0$.
The $s-1$ tree components each contain at most one further vertex.
Thus $3\le |V(W_0)|\le s+1$, and $s\ge2$.
It follows that $k=3\le s+c$.

Suppose next that $k=4$.
The core is a subdivision of $K_{3,3}$.
Delete the subdivided cycles corresponding to its nine four-cycles, one at a time.
Each deletion leaves a component containing all six branch vertices.
The complementary kernel edges form a spanning tree with two adjacent centers.
Each such component contains at most one vertex of $W_0$.

Every branch vertex belongs to all nine components.
Each kernel edge belongs to five complementary trees, since it lies in four of the nine four-cycles.
Thus an internal vertex of a subdivided edge belongs to five components.
Tree attachments again inherit the membership of their attachment vertex.
Every vertex of the cyclic component therefore belongs to at least five of the nine components.
The cyclic component contains at most $\lfloor9/5\rfloor=1$ vertex of $W_0$.
Consequently $3\le |V(W_0)|\le s$, and $s\ge3$.
Hence $k=4\le s+c$.
Equation~\ref{eq:nowaste-rank-identity} proves the bound.

For sharpness, take $G=K_5$ minus one edge.
Let $W$ be the triangle on its three vertices of degree four.
Here $F=K_{2,3}$, with the three marked vertices in its larger part.
Its outside parts are the empty set, its three four-cycles, and four paths for each of the three marked vertex pairs.
Each path closes with the corresponding edge of $W$.
The empty outside part is represented by $W$ itself.
Here $\beta(G)=5$, $\beta(W)=1$, and $|V(W)|=3$.
Thus the difference is $4=|V(W)|+1$.
\end{proof}

The simple-graph assumption supplies three vertices in $W_0$.
The theorem is not asserted for marked loops or marked parallel pairs.
The new proof has not been formalized in Lean.


\begin{corollary}[Dependence on the marked components]\label{cor:nowaste-component}
Under the same hypotheses, put $m=|V(W)|$, $c=\operatorname{comp}(W)$, and
$M=\max_i |V(W_i)|$, where $W_i$ ranges over the components of $W$.
Then
\[
 \beta(G)-\beta(W)
 \le m-c+1+\max\{0,4-M\}.
\]
In particular, the difference is at most $m$ unless $W$ is a triangle.
If $W$ is connected, the difference is at most $\max\{m,4\}$.
\end{corollary}

\begin{proof}
Keep $F,s,k$ from the preceding proof.
If $k\le1$, then $k-s\le0$, and the rank identity gives the result.
Suppose $k\ge2$.
For ranks three and four, the preceding covering arguments apply to every component $W_i$.
They bound its intersection with the cyclic component of $F$ by two and one, respectively.

For rank two, use the three cycles of the theta core.
The three complementary components contain its two branch vertices and one subdivided branch each.
Each such component contains at most one vertex of $W_i$.
Every vertex of the cyclic component belongs to at least one of them.
Thus at most three vertices of $W_i$ lie in that component.

In each of the three cases, at most $5-k$ vertices of $W_i$ lie in the cyclic component.
Each of the other $s-1$ components contributes at most one.
Consequently $|V(W_i)|\le s+4-k$, and $k-s\le4-M$.
Substitute in the rank identity.
Since $M\ge3$, an excess of one over $m$ requires $c=1$ and $M=3$.
A connected simple graph with three vertices and a cycle is a triangle.
\end{proof}

\paragraph{Equality structure.}
The equality case in Theorem~\ref{thm:nowaste-one} has an explicit description.
For a vertex in the cyclic component of $F$, call its nearest vertex in the cyclic block its \emph{anchor}.
A vertex already in the block anchors itself.
Every tree attached to the block has exactly one attachment vertex, so the anchor is well-defined.

\begin{proposition}\label{prop:nowaste-equality}
Under the no-waste hypothesis,
$\beta(G)-\beta(W)=|V(W)|+1$ if and only if $W$ is a triangle and one of the following configurations holds.
\begin{enumerate}
\item $F$ is connected of rank two.
Its cyclic block is a theta graph.
The three marked vertices have anchors in the interiors of its three different branches.
\item $F$ has two components and rank three.
Its cyclic block is a subdivision of $K_4$.
Two marked vertices lie in this component, with anchors interior to two opposite subdivided edges.
The third marked vertex lies in the tree component.
\item $F$ has three components and rank four.
Its cyclic block is a subdivision of $K_{3,3}$.
Each component of $F$ contains exactly one marked vertex.
The two other components are trees.
\end{enumerate}
Conversely, every connected simple graph with $W$ a triangle and one of these configurations satisfies the no-waste hypothesis and equality.
\end{proposition}

\begin{proof}
Corollary~\ref{cor:nowaste-component} forces $W$ to be a triangle.
The rank identity then gives $k=s+1$.
The three possibilities are $(k,s)=(2,1),(3,2),(4,3)$.

For the theta core, the three marked vertices must use the three complementary-component slots exactly once.
An anchor at a branch vertex would use all three slots by itself.
Hence the anchors are internal and lie on different branches.
For the $K_4$ core, the two marked vertices in the cyclic component must use the four star slots exactly twice each.
Thus their anchors are internal and their kernel edges have disjoint endpoints.
These edges are opposite.
For the $K_{3,3}$ core, the component bound gives at most one marked vertex in each component.
All three components therefore contain one.

For the converse, an outside part is an edge set in $F$ whose odd-degree vertices lie among the three marked vertices.
Its boundary is empty or a marked pair.
Every allowable edge set can be completed inside the triangle.
The empty-boundary sets are the even subgraphs of $F$.
Their nonempty members are cycles by the core descriptions.

In the theta configuration, each marked pair has exactly four simple paths in $F$.
These exhaust its $2^2$ joins.
In the $K_4$ configuration, only the pair in the cyclic component can be a boundary.
The eight paths between opposite-edge interior points exhaust its $2^3$ joins.
Tree tails do not change these counts or create repeated vertices.
In the $K_{3,3}$ configuration, component parity excludes every nonempty boundary.
Every nonempty outside part is therefore either a cycle of $F$ or a marked-pair path.
Close each path with the corresponding triangle edge, and use $W$ for the empty part.
This proves no waste.
The rank identity gives equality in each configuration.
\end{proof}

\paragraph{Finite checks.}
The script \texttt{check\_no\_waste.py} checks every marked edge subset containing a cycle in every connected simple graph through six vertices.
It uses the NetworkX graph atlas and exact cycle-space enumeration.
The refined component bound passes for all $101{,}904$ tested pairs, including $19{,}035$ pairs with no waste.
Every equality case has a marked triangle.
These checks support the proof but do not replace it.
